\documentclass[11pt,a4paper]{article}
\usepackage[T1]{fontenc}
\usepackage[utf8]{inputenc}
\usepackage[brazil]{babel}
\usepackage{amsmath,amssymb,textcomp}
\usepackage[margin=2.2cm]{geometry}
\usepackage{enumitem}
\usepackage{array}
\setlength{\parindent}{0pt}
\setlength{\parskip}{0.6em}

\begin{document}
\section*{Avaliação de questões --- avaliador 4}
\subsection*{Como responder}

São 22 questões de Matemática do Ensino Médio. Leia cada uma como leria uma
questão que você pensasse em usar com a sua turma, e responda às quatro
perguntas que vêm logo abaixo dela.

Onde se pergunta pelo \textbf{nível de dificuldade}, responda pensando nos
seus alunos, não num aluno ideal. É essa resposta que permite calibrar o
sistema: ele classifica cada questão como fácil, média ou difícil por conta
própria, e precisamos saber o quanto isso corresponde à sala de aula real.

\textbf{A amostra tem qualidade variável, e pode conter questões com erro
matemático.} Se encontrar algum, aponte onde.

Não há resposta certa sobre você: o que está sendo avaliado é o sistema que
produziu estas questões, não quem as lê. Um ``recusada'' bem justificado vale
mais para este trabalho do que um ``aceita'' por gentileza.

Você pode responder direto neste documento impresso ou na planilha
\texttt{respostas-avaliador-4.csv}, que já vem com os códigos dos itens nas linhas.

Tempo estimado: cerca de 75 minutos.
\bigskip\hrule\bigskip
\begin{samepage}
\subsection*{Q006}
Uma companhia de saneamento cobra a conta de água residencial de acordo com faixas de consumo mensal $m$, medido em metros cúbicos (m³), sendo o serviço válido para consumos de até $50\,m^3$. As regras de cobrança, em reais, são as seguintes:

\begin{center}
\begin{tabular}{|p{0.430\linewidth}|p{0.430\linewidth}|}
\hline
\textbf{Faixa de consumo} & \textbf{Regra de cobrança} \\
\hline
$0 \le m \le 10$ & Taxa fixa de manutenção de R\$\,5,00 mais R\$\,1,50 por m³ consumido \\
\hline
$10 < m \le 30$ & R\$\,20,00 fixos (referentes aos primeiros 10 m³) mais R\$\,2,00 por m³ que exceder 10 m³ \\
\hline
$30 < m \le 50$ & R\$\,60,00 fixos (referentes aos primeiros 30 m³) mais R\$\,4,00 por m³ que exceder 30 m³ \\
\hline
\end{tabular}
\end{center}

Seja $C(m)$ o valor da conta, em reais, para um consumo de $m$ metros cúbicos.

a) Escreva a lei algébrica de $C(m)$, definida por partes, especificando o domínio de validade de cada sentença.

b) Determine a imagem da função $C$ considerando todo o intervalo de consumo válido, de $0$ a $50\,m^3$.

c) Em certo mês, uma família pagou uma conta de R\$\,74,00. Determine o consumo desse mês, indicando claramente por que sentença da função esse valor deve ser calculado.

d) A conta de água é uma função crescente, decrescente ou constante no intervalo de consumo considerado? Justifique sua resposta levando em conta as três sentenças da função.

\textbf{Gabarito proposto:} a) $C(m) = 1{,}5m+5$ para $0\le m\le10$; $C(m)=20+2(m-10)$ para $10<m\le30$; $C(m)=60+4(m-30)$ para $30<m\le50$. b) Imagem $=[5,140]$. c) $m=33{,}5\,m^3$ (calculado na terceira sentença, pois $74>60$). d) A função é crescente em todo o domínio $[0,50]$.

\textbf{Habilidade declarada:} EM13MAT405 --- Reconhecer funções definidas por uma ou mais sentenças (como a tabela do Imposto de Renda, contas de luz, água, gás etc.), em suas representações algébrica e gráfica, convertendo essas representações de uma para outra e identificando domínios de validade, imagem, crescimento e decrescimento.
\end{samepage}

\noindent\rule{\linewidth}{0.4pt}
\textbf{Tem erro matemático?}\quad
\mbox{$\square$~não} \quad \mbox{$\square$~sim, aqui:} \underline{\hspace{5cm}}
\quad \mbox{$\square$~não sei dizer}

\textbf{Que nível de dificuldade esta questão tem para os seus alunos?}\quad
\mbox{$\square$~fácil} \quad \mbox{$\square$~média} \quad
\mbox{$\square$~difícil} \quad \mbox{$\square$~fora do alcance da turma}

\textbf{Você usaria esta questão?}\quad
\mbox{$\square$~aceita} \quad \mbox{$\square$~aceita com ajuste} \quad
\mbox{$\square$~recusada}

\textbf{Por quê?} \emph{(obrigatório quando não for ``aceita'')}

\underline{\hspace{\linewidth}}

\underline{\hspace{\linewidth}}
\vspace{0.6em}

\begin{samepage}
\subsection*{Q068}
Considere as quatro funções quadráticas a seguir, definidas para $x \in \mathbb{R}$:

$A) \; y = 3x^2$

$B) \; y = 3x^2 - 6$

$C) \; y = 3(x-2)^2$

$D) \; y = 3x^2 + 5x$

Para cada uma delas, determine algebricamente as coordenadas do vértice da parábola correspondente e, a partir dessa informação geométrica, identifique qual das quatro funções é tal que $y$ é diretamente proporcional a $x^2$ (isto é, pode ser escrita na forma $y = kx^2$). Justifique sua escolha relacionando a posição do vértice de cada parábola no plano cartesiano com a condição de proporcionalidade direta.
\begin{enumerate}[label=(\alph*),nosep,leftmargin=*]
\item A função $y = 3x^2$, pois seu gráfico é uma parábola cujo vértice está exatamente em $(0,0)$, coincidindo com a origem do plano cartesiano.
\item A função $y = 3x^2 - 6$, pois seu vértice está sobre o eixo $y$ (em $x=0$), o que já garante a proporcionalidade direta entre $y$ e $x^2$.
\item A função $y = 3(x-2)^2$, pois seu gráfico toca o eixo $x$ em um único ponto, o que caracteriza a proporcionalidade direta com $x^2$.
\item A função $y = 3x^2 + 5x$, pois seu gráfico passa pela origem (quando $x=0$, $y=0$), o que basta para garantir que $y$ é proporcional a $x^2$.
\end{enumerate}

\textbf{Gabarito proposto:} A) y = 3x², pois é a única cujo vértice está exatamente na origem (0,0) do plano cartesiano.

\textbf{Habilidade declarada:} EM13MAT402 --- Converter representações algébricas de funções polinomiais de 2º grau para representações geométricas no plano cartesiano, distinguindo os casos nos quais uma variável for diretamente proporcional ao quadrado da outra, recorrendo ou não a softwares ou aplicativos de álgebra e geometria dinâmica.
\end{samepage}

\noindent\rule{\linewidth}{0.4pt}
\textbf{Tem erro matemático?}\quad
\mbox{$\square$~não} \quad \mbox{$\square$~sim, aqui:} \underline{\hspace{5cm}}
\quad \mbox{$\square$~não sei dizer}

\textbf{Que nível de dificuldade esta questão tem para os seus alunos?}\quad
\mbox{$\square$~fácil} \quad \mbox{$\square$~média} \quad
\mbox{$\square$~difícil} \quad \mbox{$\square$~fora do alcance da turma}

\textbf{Você usaria esta questão?}\quad
\mbox{$\square$~aceita} \quad \mbox{$\square$~aceita com ajuste} \quad
\mbox{$\square$~recusada}

\textbf{Por quê?} \emph{(obrigatório quando não for ``aceita'')}

\underline{\hspace{\linewidth}}

\underline{\hspace{\linewidth}}
\vspace{0.6em}

\begin{samepage}
\subsection*{Q030}
Um engenheiro testa a queda livre de uma pequena esfera solta de um drone e mede a distância total percorrida pela esfera (em metros) para diferentes tempos de queda (em segundos), obtendo a tabela abaixo:

\begin{center}
\begin{tabular}{|p{0.430\linewidth}|p{0.430\linewidth}|}
\hline
\textbf{Tempo $t$ (s)} & \textbf{Distância $d$ (m)} \\
\hline
1 & 5 \\
\hline
2 & 20 \\
\hline
3 & 45 \\
\hline
4 & 80 \\
\hline
\end{tabular}
\end{center}

a) Analisando os valores da tabela, identifique o padrão que relaciona $d$ e $t$ e escreva a lei algébrica que expressa $d$ em função de $t$.

b) Usando essa lei, determine a distância percorrida pela esfera após 5 segundos de queda.

\textbf{Gabarito proposto:} a) $d(t) = 5t^2$ (a distância é proporcional ao quadrado do tempo). b) $d(5) = 125$ m.

\textbf{Habilidade declarada:} EM13MAT502 --- Investigar relações entre números expressos em tabelas para representá-los no plano cartesiano, identificando padrões e criando conjecturas para generalizar e expressar algebricamente essa generalização, reconhecendo quando essa representação é de função polinomial de 2º grau do tipo y = ax².
\end{samepage}

\noindent\rule{\linewidth}{0.4pt}
\textbf{Tem erro matemático?}\quad
\mbox{$\square$~não} \quad \mbox{$\square$~sim, aqui:} \underline{\hspace{5cm}}
\quad \mbox{$\square$~não sei dizer}

\textbf{Que nível de dificuldade esta questão tem para os seus alunos?}\quad
\mbox{$\square$~fácil} \quad \mbox{$\square$~média} \quad
\mbox{$\square$~difícil} \quad \mbox{$\square$~fora do alcance da turma}

\textbf{Você usaria esta questão?}\quad
\mbox{$\square$~aceita} \quad \mbox{$\square$~aceita com ajuste} \quad
\mbox{$\square$~recusada}

\textbf{Por quê?} \emph{(obrigatório quando não for ``aceita'')}

\underline{\hspace{\linewidth}}

\underline{\hspace{\linewidth}}
\vspace{0.6em}

\begin{samepage}
\subsection*{Q019}
Um biólogo estuda o crescimento de uma cultura de bactérias em laboratório. O número de bactérias, em milhares, após $t$ horas (com $t$ podendo assumir qualquer valor real, inclusive negativo, para fins de modelagem matemática) é dado por $f(t) = 2^t$. Para descobrir quanto tempo é necessário para que a cultura atinja uma quantidade $N$ (em milhares) de bactérias, o biólogo utiliza a função inversa $g(N) = \log_2(N)$. As duas funções são representadas no mesmo plano cartesiano. Com base na análise do domínio, da imagem e do comportamento de crescimento de $f$ e $g$, assinale a alternativa correta.
\begin{enumerate}[label=(\alph*),nosep,leftmargin=*]
\item O domínio de f (todos os reais) é igual à imagem de g, e a imagem de f (reais positivos) é igual ao domínio de g; ambas as funções são crescentes em seus domínios.
\item O domínio de f é o conjunto dos reais positivos e a imagem de g é o conjunto de todos os reais, sendo f decrescente e g crescente.
\item O domínio de f coincide com a imagem de g, mas g é decrescente enquanto f é crescente.
\item Como f e g são inversas, ambas têm exatamente o mesmo domínio (os reais) e a mesma imagem (os reais positivos), sendo ambas crescentes.
\end{enumerate}

\textbf{Gabarito proposto:} O domínio de f (todos os reais) é igual à imagem de g, e a imagem de f (reais positivos) é igual ao domínio de g; ambas as funções são crescentes em seus domínios.

\textbf{Habilidade declarada:} EM13MAT403 --- Comparar e analisar as representações, em plano cartesiano, das funções exponencial e logarítmica para identificar as características fundamentais (domínio, imagem, crescimento) de cada uma, com ou sem apoio de tecnologias digitais, estabelecendo relações entre elas.
\end{samepage}

\noindent\rule{\linewidth}{0.4pt}
\textbf{Tem erro matemático?}\quad
\mbox{$\square$~não} \quad \mbox{$\square$~sim, aqui:} \underline{\hspace{5cm}}
\quad \mbox{$\square$~não sei dizer}

\textbf{Que nível de dificuldade esta questão tem para os seus alunos?}\quad
\mbox{$\square$~fácil} \quad \mbox{$\square$~média} \quad
\mbox{$\square$~difícil} \quad \mbox{$\square$~fora do alcance da turma}

\textbf{Você usaria esta questão?}\quad
\mbox{$\square$~aceita} \quad \mbox{$\square$~aceita com ajuste} \quad
\mbox{$\square$~recusada}

\textbf{Por quê?} \emph{(obrigatório quando não for ``aceita'')}

\underline{\hspace{\linewidth}}

\underline{\hspace{\linewidth}}
\vspace{0.6em}

\begin{samepage}
\subsection*{Q024}
A tabela abaixo relaciona valores de $x$ e os correspondentes valores de $y$, obtidos a partir de uma mesma regra de formação (os valores de $x$ não variam em intervalos iguais):

\begin{center}
\begin{tabular}{|p{0.172\linewidth}|p{0.172\linewidth}|p{0.172\linewidth}|p{0.172\linewidth}|p{0.172\linewidth}|}
\hline
\textbf{$x$} & \textbf{$-3$} & \textbf{$1$} & \textbf{$4$} & \textbf{$8$} \\
\hline
$y$ & $14$ & $2$ & $-7$ & $-19$ \\
\hline
\end{tabular}
\end{center}

Um estudante, analisando os pares ordenados $(x,y)$ da tabela, conjectura que essa relação pode ser representada por uma função polinomial do 1º grau. Assinale a alternativa que apresenta a lei de formação $f(x)$ que generaliza corretamente o padrão observado na tabela.
\begin{enumerate}[label=(\alph*),nosep,leftmargin=*]
\item $f(x) = -3x + 5$
\item $f(x) = -12x + 14$
\item $f(x) = 3x - 1$
\item $f(x) = -11x + 13$
\end{enumerate}

\textbf{Gabarito proposto:} f(x) = -3x + 5

\textbf{Habilidade declarada:} EM13MAT501 --- Investigar relações entre números expressos em tabelas para representá-los no plano cartesiano, identificando padrões e criando conjecturas para generalizar e expressar algebricamente essa generalização, reconhecendo quando essa representação é de função polinomial de 1º grau.
\end{samepage}

\noindent\rule{\linewidth}{0.4pt}
\textbf{Tem erro matemático?}\quad
\mbox{$\square$~não} \quad \mbox{$\square$~sim, aqui:} \underline{\hspace{5cm}}
\quad \mbox{$\square$~não sei dizer}

\textbf{Que nível de dificuldade esta questão tem para os seus alunos?}\quad
\mbox{$\square$~fácil} \quad \mbox{$\square$~média} \quad
\mbox{$\square$~difícil} \quad \mbox{$\square$~fora do alcance da turma}

\textbf{Você usaria esta questão?}\quad
\mbox{$\square$~aceita} \quad \mbox{$\square$~aceita com ajuste} \quad
\mbox{$\square$~recusada}

\textbf{Por quê?} \emph{(obrigatório quando não for ``aceita'')}

\underline{\hspace{\linewidth}}

\underline{\hspace{\linewidth}}
\vspace{0.6em}

\begin{samepage}
\subsection*{Q015}
Uma colônia de bactérias em um experimento de laboratório tem seu número de indivíduos, a partir do início da observação, dado por $N(t) = 100 \cdot 2^{t}$, em que $t$ é o tempo em horas ($t \geq 0$) e $N(t)$ é o número de bactérias. Um pesquisador quer saber, ao contrário, quanto tempo é necessário para que a colônia atinja um determinado número $N$ de bactérias, e por isso utiliza a função inversa de $N(t)$, que ele chama de $t(N)$.

a) Determine o domínio e a imagem de $N(t)$, considerando o contexto do experimento.

b) Encontre a expressão de $t(N)$ (a função inversa de $N(t)$) e determine seu domínio e sua imagem.

c) Compare o crescimento das duas funções: as duas são crescentes, mas uma cresce 'cada vez mais rápido' e a outra cresce 'cada vez mais devagar' à medida que a variável aumenta. Identifique qual é qual e justifique observando como varia o incremento de cada função.

d) Explique por que o domínio de $N(t)$ coincide com a imagem de $t(N)$, e a imagem de $N(t)$ coincide com o domínio de $t(N)$, relacionando esse fato com a posição dos gráficos das duas funções no plano cartesiano.

\textbf{Gabarito proposto:} a) $D_N=[0,+\infty)$, $Im_N=[100,+\infty)$. b) $t(N)=\log_2(N/100)$, com $D_t=[100,+\infty)$, $Im_t=[0,+\infty)$. c) $N(t)$ cresce cada vez mais rápido (convexa); $t(N)$ cresce cada vez mais devagar (côncava). d) Por serem funções inversas, seus gráficos são simétricos em relação à reta $y=x$, o que troca domínio e imagem entre as duas funções.

\textbf{Habilidade declarada:} EM13MAT403 --- Comparar e analisar as representações, em plano cartesiano, das funções exponencial e logarítmica para identificar as características fundamentais (domínio, imagem, crescimento) de cada uma, com ou sem apoio de tecnologias digitais, estabelecendo relações entre elas.
\end{samepage}

\noindent\rule{\linewidth}{0.4pt}
\textbf{Tem erro matemático?}\quad
\mbox{$\square$~não} \quad \mbox{$\square$~sim, aqui:} \underline{\hspace{5cm}}
\quad \mbox{$\square$~não sei dizer}

\textbf{Que nível de dificuldade esta questão tem para os seus alunos?}\quad
\mbox{$\square$~fácil} \quad \mbox{$\square$~média} \quad
\mbox{$\square$~difícil} \quad \mbox{$\square$~fora do alcance da turma}

\textbf{Você usaria esta questão?}\quad
\mbox{$\square$~aceita} \quad \mbox{$\square$~aceita com ajuste} \quad
\mbox{$\square$~recusada}

\textbf{Por quê?} \emph{(obrigatório quando não for ``aceita'')}

\underline{\hspace{\linewidth}}

\underline{\hspace{\linewidth}}
\vspace{0.6em}

\begin{samepage}
\subsection*{Q041}
Um ponto P parte da posição (0, 5) no instante t = 0 e passa a se mover indefinidamente, no sentido anti-horário, sobre uma circunferência de centro (0, 3) e raio 2, completando uma volta completa a cada intervalo de tempo $\Delta$t = 4$\pi$ unidades de tempo. A ordenada de P, em função do tempo t, é dada por y(t) = 3 + 2cos(t/2).

Comparando a trajetória de P na circunferência (o ciclo trigonométrico deslocado e ampliado) com a representação gráfica de y(t) no plano cartesiano t × y, assinale a alternativa que apresenta corretamente o domínio, a imagem e o período dessa função.
\begin{enumerate}[label=(\alph*),nosep,leftmargin=*]
\item Domínio = $\mathbb{R}$; Imagem = [1, 5]; Período = 4$\pi$
\item Domínio = [0, +$\infty$); Imagem = [-2, 2]; Período = 4$\pi$
\item Domínio = [0, +$\infty$); Imagem = [1, 5]; Período = 2$\pi$
\item Domínio = [0, +$\infty$); Imagem = [1, 5]; Período = 4$\pi$
\end{enumerate}

\textbf{Gabarito proposto:} D) Domínio = [0, +$\infty$); Imagem = [1, 5]; Período = 4$\pi$

\textbf{Habilidade declarada:} EM13MAT404 --- Identificar as características fundamentais das funções seno e cosseno (periodicidade, domínio, imagem), por meio da comparação das representações em ciclos trigonométricos e em planos cartesianos, com ou sem apoio de tecnologias digitais.
\end{samepage}

\noindent\rule{\linewidth}{0.4pt}
\textbf{Tem erro matemático?}\quad
\mbox{$\square$~não} \quad \mbox{$\square$~sim, aqui:} \underline{\hspace{5cm}}
\quad \mbox{$\square$~não sei dizer}

\textbf{Que nível de dificuldade esta questão tem para os seus alunos?}\quad
\mbox{$\square$~fácil} \quad \mbox{$\square$~média} \quad
\mbox{$\square$~difícil} \quad \mbox{$\square$~fora do alcance da turma}

\textbf{Você usaria esta questão?}\quad
\mbox{$\square$~aceita} \quad \mbox{$\square$~aceita com ajuste} \quad
\mbox{$\square$~recusada}

\textbf{Por quê?} \emph{(obrigatório quando não for ``aceita'')}

\underline{\hspace{\linewidth}}

\underline{\hspace{\linewidth}}
\vspace{0.6em}

\begin{samepage}
\subsection*{Q070}
A magnitude de um abalo sísmico na escala Richter é definida por $M = \log_{10}\left(\dfrac{I}{I_0}\right)$, em que $I$ é a intensidade da onda sísmica registrada pelo sismógrafo e $I_0$ é uma intensidade de referência fixa, igual para todos os abalos.

Um sismólogo estuda uma situação hipotética em que dois abalos, de magnitudes $M_1$ e $M_2$, ocorrem na mesma falha geológica e suas ondas se sobrepõem exatamente no mesmo instante, de modo que a intensidade total registrada pelo sismógrafo passa a ser a soma das intensidades individuais dos dois abalos, isto é, $I = I_1 + I_2$.

Elabore a expressão geral que forneça a magnitude resultante $M$ em função de $M_1$ e $M_2$ nessa situação e utilize-a para calcular o valor de $M$ quando $M_1 = 5$ e $M_2 = 6$.

Assinale a alternativa que apresenta corretamente o valor de $M$ (aproximado a duas casas decimais).
\begin{enumerate}[label=(\alph*),nosep,leftmargin=*]
\item $M \approx 6{,}04$
\item $M = 11$
\item $M = 5{,}5$
\item $M = 6$
\end{enumerate}

\textbf{Gabarito proposto:} M $\approx$ 6,04

\textbf{Habilidade declarada:} EM13MAT305 --- Resolver e elaborar problemas com funções logarítmicas nos quais é necessário compreender e interpretar a variação das grandezas envolvidas, em contextos como os de abalos sísmicos, pH, radioatividade, Matemática Financeira, entre outros.
\end{samepage}

\noindent\rule{\linewidth}{0.4pt}
\textbf{Tem erro matemático?}\quad
\mbox{$\square$~não} \quad \mbox{$\square$~sim, aqui:} \underline{\hspace{5cm}}
\quad \mbox{$\square$~não sei dizer}

\textbf{Que nível de dificuldade esta questão tem para os seus alunos?}\quad
\mbox{$\square$~fácil} \quad \mbox{$\square$~média} \quad
\mbox{$\square$~difícil} \quad \mbox{$\square$~fora do alcance da turma}

\textbf{Você usaria esta questão?}\quad
\mbox{$\square$~aceita} \quad \mbox{$\square$~aceita com ajuste} \quad
\mbox{$\square$~recusada}

\textbf{Por quê?} \emph{(obrigatório quando não for ``aceita'')}

\underline{\hspace{\linewidth}}

\underline{\hspace{\linewidth}}
\vspace{0.6em}

\begin{samepage}
\subsection*{Q058}
Uma agência de viagens vende pacotes turísticos a um preço unitário de $x$ reais. Um estudo de mercado mostrou que a quantidade de pacotes vendidos por mês, em função do preço, é dada por $q(x) = 1000 - 4x$ (em unidades). Por cláusula do contrato firmado com a operadora parceira, o preço do pacote deve obrigatoriamente estar entre R$\,200,00$ e R$\,300,00$, incluindo os extremos. Considerando a receita mensal $R(x) = x \cdot q(x)$, determine o preço que maximiza a receita da agência dentro dessa faixa contratual e o valor dessa receita máxima.
\begin{enumerate}[label=(\alph*),nosep,leftmargin=*]
\item R\$\,125,00, com receita máxima de R\$\,62.500,00
\item R\$\,200,00, com receita máxima de R\$\,40.000,00
\item R\$\,300,00, com receita máxima de R\$-60.000,00
\item R\$\,250,00, com receita máxima de R\$\,0,00
\end{enumerate}

\textbf{Gabarito proposto:} Preço de R\$\,200,00, com receita máxima de R\$\,40.000,00 (alternativa B)

\textbf{Habilidade declarada:} EM13MAT503 --- Investigar pontos de máximo ou de mínimo de funções quadráticas em contextos da Matemática Financeira ou da Cinemática, entre outros.
\end{samepage}

\noindent\rule{\linewidth}{0.4pt}
\textbf{Tem erro matemático?}\quad
\mbox{$\square$~não} \quad \mbox{$\square$~sim, aqui:} \underline{\hspace{5cm}}
\quad \mbox{$\square$~não sei dizer}

\textbf{Que nível de dificuldade esta questão tem para os seus alunos?}\quad
\mbox{$\square$~fácil} \quad \mbox{$\square$~média} \quad
\mbox{$\square$~difícil} \quad \mbox{$\square$~fora do alcance da turma}

\textbf{Você usaria esta questão?}\quad
\mbox{$\square$~aceita} \quad \mbox{$\square$~aceita com ajuste} \quad
\mbox{$\square$~recusada}

\textbf{Por quê?} \emph{(obrigatório quando não for ``aceita'')}

\underline{\hspace{\linewidth}}

\underline{\hspace{\linewidth}}
\vspace{0.6em}

\begin{samepage}
\subsection*{Q007}
Uma oficina mecânica cobra por seus serviços de revisão um valor fixo de mão de obra, mais uma taxa que depende do número de horas de trabalho. Uma revisão que durou 2 horas custou R\$\,150,00, e outra revisão, no mesmo padrão de cobrança, durou 5 horas e custou R\$\,300,00. Sabendo que o custo total varia linearmente com o número de horas trabalhadas, qual seria o custo de uma revisão que durasse 8 horas?
\begin{enumerate}[label=(\alph*),nosep,leftmargin=*]
\item R\$\,450,00
\item R\$\,400,00
\item R\$\,600,00
\item R\$\,480,00
\end{enumerate}

\textbf{Gabarito proposto:} R\$\,450,00

\textbf{Habilidade declarada:} EM13MAT302 --- Resolver e elaborar problemas cujos modelos são as funções polinomiais de 1º e 2º graus, em contextos diversos, incluindo ou não tecnologias digitais.
\end{samepage}

\noindent\rule{\linewidth}{0.4pt}
\textbf{Tem erro matemático?}\quad
\mbox{$\square$~não} \quad \mbox{$\square$~sim, aqui:} \underline{\hspace{5cm}}
\quad \mbox{$\square$~não sei dizer}

\textbf{Que nível de dificuldade esta questão tem para os seus alunos?}\quad
\mbox{$\square$~fácil} \quad \mbox{$\square$~média} \quad
\mbox{$\square$~difícil} \quad \mbox{$\square$~fora do alcance da turma}

\textbf{Você usaria esta questão?}\quad
\mbox{$\square$~aceita} \quad \mbox{$\square$~aceita com ajuste} \quad
\mbox{$\square$~recusada}

\textbf{Por quê?} \emph{(obrigatório quando não for ``aceita'')}

\underline{\hspace{\linewidth}}

\underline{\hspace{\linewidth}}
\vspace{0.6em}

\begin{samepage}
\subsection*{Q066}
Duas progressões aritméticas de números naturais são dadas por:

$a_n = 4n - 1,\; n \in \mathbb{N}^*$ (ou seja, $3, 7, 11, 15, 19, 23, 27, \dots$)

$b_m = 6m - 3,\; m \in \mathbb{N}^*$ (ou seja, $3, 9, 15, 21, 27, \dots$)

Cada uma dessas progressões pode ser vista como a restrição ao conjunto dos números naturais positivos de uma função afim: $f(n) = 4n-1$ e $g(m) = 6m-3$, respectivamente, com $n, m$ percorrendo $\mathbb{N}^*$.

Alguns números aparecem nas duas listas (por exemplo, $3$, $15$ e $27$). Chame de $c_k$, com $k = 0, 1, 2, \dots$, o $(k+1)$-ésimo desses termos comuns, em ordem crescente.

a) Mostre que a sequência $(c_k)$ também é uma progressão aritmética e determine seu primeiro termo $c_0$ e sua razão.

b) Escreva a função afim $h: \mathbb{N} \to \mathbb{R}$ tal que $h(k) = c_k$ para todo $k \in \mathbb{N}$.

c) Quantos termos comuns às duas progressões existem entre $1$ e $2023$, inclusive?

\textbf{Gabarito proposto:} a) $(c_k)$ é PA com $c_0=3$ e razão $12$. b) $h(k)=12k+3$, $k\in\mathbb{N}$. c) $169$ termos comuns.

\textbf{Habilidade declarada:} EM13MAT507 --- Identificar e associar sequências numéricas (PA) a funções afins de domínios discretos para análise de propriedades, incluindo dedução de algumas fórmulas e resolução de problemas.
\end{samepage}

\noindent\rule{\linewidth}{0.4pt}
\textbf{Tem erro matemático?}\quad
\mbox{$\square$~não} \quad \mbox{$\square$~sim, aqui:} \underline{\hspace{5cm}}
\quad \mbox{$\square$~não sei dizer}

\textbf{Que nível de dificuldade esta questão tem para os seus alunos?}\quad
\mbox{$\square$~fácil} \quad \mbox{$\square$~média} \quad
\mbox{$\square$~difícil} \quad \mbox{$\square$~fora do alcance da turma}

\textbf{Você usaria esta questão?}\quad
\mbox{$\square$~aceita} \quad \mbox{$\square$~aceita com ajuste} \quad
\mbox{$\square$~recusada}

\textbf{Por quê?} \emph{(obrigatório quando não for ``aceita'')}

\underline{\hspace{\linewidth}}

\underline{\hspace{\linewidth}}
\vspace{0.6em}

\begin{samepage}
\subsection*{Q087}
Considere a progressão aritmética $(a_n)$ com primeiro termo $a_1 = 7$ e razão $r = 4$: $(7, 11, 15, 19, 23, \dots)$. Essa PA pode ser vista como a restrição da função afim $f(x) = 4x + 3$ ao conjunto dos números naturais não nulos, de modo que $a_n = f(n)$ para todo $n \geq 1$ (o gráfico da PA é o conjunto de pontos discretos $(n, a_n)$ sobre a reta que representa $f$).

Qual é o menor termo dessa PA que é maior do que 100?
\begin{enumerate}[label=(\alph*),nosep,leftmargin=*]
\item 103
\item 99
\item 107
\item 100
\end{enumerate}

\textbf{Gabarito proposto:} 103

\textbf{Habilidade declarada:} EM13MAT507 --- Identificar e associar sequências numéricas (PA) a funções afins de domínios discretos para análise de propriedades, incluindo dedução de algumas fórmulas e resolução de problemas.
\end{samepage}

\noindent\rule{\linewidth}{0.4pt}
\textbf{Tem erro matemático?}\quad
\mbox{$\square$~não} \quad \mbox{$\square$~sim, aqui:} \underline{\hspace{5cm}}
\quad \mbox{$\square$~não sei dizer}

\textbf{Que nível de dificuldade esta questão tem para os seus alunos?}\quad
\mbox{$\square$~fácil} \quad \mbox{$\square$~média} \quad
\mbox{$\square$~difícil} \quad \mbox{$\square$~fora do alcance da turma}

\textbf{Você usaria esta questão?}\quad
\mbox{$\square$~aceita} \quad \mbox{$\square$~aceita com ajuste} \quad
\mbox{$\square$~recusada}

\textbf{Por quê?} \emph{(obrigatório quando não for ``aceita'')}

\underline{\hspace{\linewidth}}

\underline{\hspace{\linewidth}}
\vspace{0.6em}

\begin{samepage}
\subsection*{Q065}
Um oceanógrafo registra a altura da maré (em metros) em um pequeno porto ao longo de um dia. Ele observa que a maré varia de forma periódica e simétrica em torno de um nível médio, atingindo seu valor máximo de 1,8 m exatamente à meia-noite ($t=0$, com $t$ em horas) e seu valor mínimo de 0,2 m exatamente às 6h da manhã. Esse padrão se repete a cada 12 horas.

a) Modele a altura da maré, em metros, por uma função do tipo $h(t) = A\cos(Bt) + D$, com $t$ em horas contado a partir da meia-noite. Determine explicitamente os valores de $A$, $B$ e $D$, justificando cada um a partir dos dados do fenômeno (amplitude, período e deslocamento vertical do gráfico em relação ao eixo $t$).

b) Determine todos os horários dentro das primeiras 12 horas do dia (isto é, para $0 \le t < 12$) em que a maré atinge exatamente 1,4 m de altura. Para cada horário encontrado, indique se a maré está subindo ou descendo naquele instante, justificando sua resposta apenas a partir do comportamento gráfico da função cosseno (sem usar derivadas).

\textbf{Gabarito proposto:} h(t) = 0,8·cos($\pi$t/6) + 1,0 (A=0,8 m; B=$\pi$/6 rad/h; D=1,0 m). A maré atinge 1,4 m às 2h (descendo) e às 10h (subindo).

\textbf{Habilidade declarada:} EM13MAT306 --- Resolver e elaborar problemas em contextos que envolvem fenômenos periódicos reais, como ondas sonoras, ciclos menstruais, movimentos cíclicos, entre outros, e comparar suas representações com as funções seno e cosseno, no plano cartesiano, com ou sem apoio de aplicativos de álgebra e geometria.
\end{samepage}

\noindent\rule{\linewidth}{0.4pt}
\textbf{Tem erro matemático?}\quad
\mbox{$\square$~não} \quad \mbox{$\square$~sim, aqui:} \underline{\hspace{5cm}}
\quad \mbox{$\square$~não sei dizer}

\textbf{Que nível de dificuldade esta questão tem para os seus alunos?}\quad
\mbox{$\square$~fácil} \quad \mbox{$\square$~média} \quad
\mbox{$\square$~difícil} \quad \mbox{$\square$~fora do alcance da turma}

\textbf{Você usaria esta questão?}\quad
\mbox{$\square$~aceita} \quad \mbox{$\square$~aceita com ajuste} \quad
\mbox{$\square$~recusada}

\textbf{Por quê?} \emph{(obrigatório quando não for ``aceita'')}

\underline{\hspace{\linewidth}}

\underline{\hspace{\linewidth}}
\vspace{0.6em}

\begin{samepage}
\subsection*{Q023}
Um biólogo estuda o crescimento de uma cultura de bactérias em laboratório. A quantidade de bactérias na cultura, em milhares, $t$ horas após o início da observação, é dada por $N(t) = 100 \cdot 2^{t/3}$, para $t \geq 0$. Para planejar a coleta de amostras, o biólogo também utiliza uma função que fornece o tempo necessário, em horas, para que a cultura atinja uma quantidade $N$ de bactérias (em milhares): $T(N) = 3 \cdot \log_2\left(\dfrac{N}{100}\right)$, para $N \geq 100$. Analisando o domínio, a imagem e o comportamento de crescimento das duas funções, bem como a relação entre seus gráficos, assinale a alternativa correta.
\begin{enumerate}[label=(\alph*),nosep,leftmargin=*]
\item O domínio de $N(t)$ é $[0,+\infty)$ e sua imagem é $[100,+\infty)$; o domínio de $T(N)$ é $[100,+\infty)$ e sua imagem é $[0,+\infty)$. Ambas as funções são crescentes, e seus gráficos são simétricos em relação à reta $y=x$, pois $T$ é a função inversa de $N$.
\item O domínio de $N(t)$ é $[100,+\infty)$ e sua imagem é $[0,+\infty)$; o domínio de $T(N)$ é $[0,+\infty)$ e sua imagem é $[100,+\infty)$. Ambas as funções são crescentes, e seus gráficos são simétricos em relação à reta $y=x$.
\item O domínio de $N(t)$ é $[0,+\infty)$ e sua imagem é $[100,+\infty)$; o domínio de $T(N)$ é $[100,+\infty)$ e sua imagem é $[0,+\infty)$. A função $N(t)$ é crescente, mas $T(N)$ é decrescente, pois o logaritmo cresce mais lentamente que a exponencial.
\item O domínio de $N(t)$ é $[0,+\infty)$ e sua imagem é $[100,+\infty)$; como toda função logarítmica está definida para qualquer número real, o domínio de $T(N)$ é $\mathbb{R}$, e sua imagem é $[0,+\infty)$. Ambas as funções são crescentes.
\end{enumerate}

\textbf{Gabarito proposto:} Domínio de $N(t)$: $[0,+\infty)$; Imagem de $N(t)$: $[100,+\infty)$; Domínio de $T(N)$: $[100,+\infty)$; Imagem de $T(N)$: $[0,+\infty)$; ambas crescentes; gráficos simétricos em relação a $y=x$ (funções inversas).

\textbf{Habilidade declarada:} EM13MAT403 --- Comparar e analisar as representações, em plano cartesiano, das funções exponencial e logarítmica para identificar as características fundamentais (domínio, imagem, crescimento) de cada uma, com ou sem apoio de tecnologias digitais, estabelecendo relações entre elas.
\end{samepage}

\noindent\rule{\linewidth}{0.4pt}
\textbf{Tem erro matemático?}\quad
\mbox{$\square$~não} \quad \mbox{$\square$~sim, aqui:} \underline{\hspace{5cm}}
\quad \mbox{$\square$~não sei dizer}

\textbf{Que nível de dificuldade esta questão tem para os seus alunos?}\quad
\mbox{$\square$~fácil} \quad \mbox{$\square$~média} \quad
\mbox{$\square$~difícil} \quad \mbox{$\square$~fora do alcance da turma}

\textbf{Você usaria esta questão?}\quad
\mbox{$\square$~aceita} \quad \mbox{$\square$~aceita com ajuste} \quad
\mbox{$\square$~recusada}

\textbf{Por quê?} \emph{(obrigatório quando não for ``aceita'')}

\underline{\hspace{\linewidth}}

\underline{\hspace{\linewidth}}
\vspace{0.6em}

\begin{samepage}
\subsection*{Q089}
Uma função quadrática $y = ax^2 + bx + c$ (com $a \neq 0$) representa uma relação em que $y$ é diretamente proporcional ao quadrado de $x$ quando existe uma constante $k \neq 0$ tal que $y = kx^2$ para todo $x$ real — ou seja, quando $b = 0$ e $c = 0$ simultaneamente.

Considere quatro parábolas, todas com concavidade voltada para cima, esboçadas no plano cartesiano e descritas apenas por características geométricas de seus gráficos:

\begin{itemize}[nosep,leftmargin=*]
\item \textbf{Parábola A}: intercepta o eixo $x$ nos pontos $(0,0)$ e $(4,0)$, e passa pelo ponto $(1,-3)$.
\item \textbf{Parábola B}: tangencia o eixo $x$ exatamente no ponto $(0,0)$ — esse ponto é o vértice da parábola — e passa pelo ponto $(3,18)$.
\item \textbf{Parábola C}: intercepta o eixo $x$ nos pontos $(-3,0)$ e $(3,0)$, e passa pelo ponto $(1,-8)$.
\item \textbf{Parábola D}: tangencia o eixo $x$ exatamente no ponto $(2,0)$ — esse ponto é o vértice da parábola — e passa pelo ponto $(0,8)$.
\end{itemize}

Assinale a alternativa que identifica corretamente qual das parábolas descritas representa uma função em que $y$ é diretamente proporcional ao quadrado de $x$.
\begin{enumerate}[label=(\alph*),nosep,leftmargin=*]
\item Parábola A, pois seu gráfico passa pela origem do plano cartesiano.
\item Parábola B, pois seu vértice está na origem e ela não possui termo linear nem constante.
\item Parábola C, pois seu gráfico é simétrico em relação ao eixo y.
\item Parábola D, pois seu gráfico tangencia o eixo x em um único ponto.
\end{enumerate}

\textbf{Gabarito proposto:} Alternativa (b): Parábola B, pois $y=2x^2$, com vértice na origem, $b=0$ e $c=0$.

\textbf{Habilidade declarada:} EM13MAT402 --- Converter representações algébricas de funções polinomiais de 2º grau para representações geométricas no plano cartesiano, distinguindo os casos nos quais uma variável for diretamente proporcional ao quadrado da outra, recorrendo ou não a softwares ou aplicativos de álgebra e geometria dinâmica.
\end{samepage}

\noindent\rule{\linewidth}{0.4pt}
\textbf{Tem erro matemático?}\quad
\mbox{$\square$~não} \quad \mbox{$\square$~sim, aqui:} \underline{\hspace{5cm}}
\quad \mbox{$\square$~não sei dizer}

\textbf{Que nível de dificuldade esta questão tem para os seus alunos?}\quad
\mbox{$\square$~fácil} \quad \mbox{$\square$~média} \quad
\mbox{$\square$~difícil} \quad \mbox{$\square$~fora do alcance da turma}

\textbf{Você usaria esta questão?}\quad
\mbox{$\square$~aceita} \quad \mbox{$\square$~aceita com ajuste} \quad
\mbox{$\square$~recusada}

\textbf{Por quê?} \emph{(obrigatório quando não for ``aceita'')}

\underline{\hspace{\linewidth}}

\underline{\hspace{\linewidth}}
\vspace{0.6em}

\begin{samepage}
\subsection*{Q018}
Uma empresa de saneamento básico cobra mensalmente pelo consumo de água uma tarifa composta por um valor fixo somado a um valor proporcional ao consumo, medido em metros cúbicos (m³). A tabela a seguir mostra o valor total pago por três clientes que tiveram consumos diferentes num mesmo mês:

\begin{center}
\begin{tabular}{|p{0.430\linewidth}|p{0.430\linewidth}|}
\hline
\textbf{Consumo (m³)} & \textbf{Valor cobrado (R\$)} \\
\hline
3 & 21,50 \\
\hline
7 & 37,50 \\
\hline
12 & 57,50 \\
\hline
\end{tabular}
\end{center}

a) Analise os dados da tabela e explique por que o valor cobrado V pode ser descrito por uma função polinomial do 1º grau do consumo x.

b) Determine a lei algébrica que expressa V em função de x.

c) Qual é o valor da tarifa fixa mensal, isto é, a parcela que o cliente paga mesmo que não haja variação proporcional ao consumo?

\textbf{Gabarito proposto:} A relação é uma função afim porque a razão $\Delta V/\Delta x$ é constante e igual a 4, mesmo com intervalos de consumo diferentes. A lei é $V(x) = 4x + 9{,}5$, e a tarifa fixa é R\$\,9,50.

\textbf{Habilidade declarada:} EM13MAT501 --- Investigar relações entre números expressos em tabelas para representá-los no plano cartesiano, identificando padrões e criando conjecturas para generalizar e expressar algebricamente essa generalização, reconhecendo quando essa representação é de função polinomial de 1º grau.
\end{samepage}

\noindent\rule{\linewidth}{0.4pt}
\textbf{Tem erro matemático?}\quad
\mbox{$\square$~não} \quad \mbox{$\square$~sim, aqui:} \underline{\hspace{5cm}}
\quad \mbox{$\square$~não sei dizer}

\textbf{Que nível de dificuldade esta questão tem para os seus alunos?}\quad
\mbox{$\square$~fácil} \quad \mbox{$\square$~média} \quad
\mbox{$\square$~difícil} \quad \mbox{$\square$~fora do alcance da turma}

\textbf{Você usaria esta questão?}\quad
\mbox{$\square$~aceita} \quad \mbox{$\square$~aceita com ajuste} \quad
\mbox{$\square$~recusada}

\textbf{Por quê?} \emph{(obrigatório quando não for ``aceita'')}

\underline{\hspace{\linewidth}}

\underline{\hspace{\linewidth}}
\vspace{0.6em}

\begin{samepage}
\subsection*{Q079}
Uma marcenaria iniciou a fabricação de cadeiras artesanais em regime de produção mensal constante. No primeiro mês de operação ($n=1$) foram fabricadas 120 cadeiras. A partir do segundo mês, a quantidade produzida a cada mês aumenta sempre do mesmo valor fixo em relação ao mês anterior, de modo que no quarto mês ($n=4$) a produção atingiu 165 cadeiras. Sabe-se que esse ritmo de crescimento se mantém constante indefinidamente, e que $n$ representa o número do mês de produção ($n = 1, 2, 3, \dots$).

a) Determine a razão da progressão aritmética formada pela produção mensal e escreva a lei de uma função afim $f(n)$ que forneça a quantidade de cadeiras fabricadas no mês $n$, indicando explicitamente o domínio dessa função.

b) Usando a função $f(n)$ obtida, determine em que mês a produção mensal será de 300 cadeiras.

c) Calcule o total de cadeiras fabricadas do 1º ao 10º mês.

\textbf{Gabarito proposto:} a) $r = 15$; $f(n) = 15n + 105$, com $n \in \mathbb{N}^* = \{1,2,3,\dots\}$. b) 13º mês. c) 1875 cadeiras.

\textbf{Habilidade declarada:} EM13MAT507 --- Identificar e associar sequências numéricas (PA) a funções afins de domínios discretos para análise de propriedades, incluindo dedução de algumas fórmulas e resolução de problemas.
\end{samepage}

\noindent\rule{\linewidth}{0.4pt}
\textbf{Tem erro matemático?}\quad
\mbox{$\square$~não} \quad \mbox{$\square$~sim, aqui:} \underline{\hspace{5cm}}
\quad \mbox{$\square$~não sei dizer}

\textbf{Que nível de dificuldade esta questão tem para os seus alunos?}\quad
\mbox{$\square$~fácil} \quad \mbox{$\square$~média} \quad
\mbox{$\square$~difícil} \quad \mbox{$\square$~fora do alcance da turma}

\textbf{Você usaria esta questão?}\quad
\mbox{$\square$~aceita} \quad \mbox{$\square$~aceita com ajuste} \quad
\mbox{$\square$~recusada}

\textbf{Por quê?} \emph{(obrigatório quando não for ``aceita'')}

\underline{\hspace{\linewidth}}

\underline{\hspace{\linewidth}}
\vspace{0.6em}

\begin{samepage}
\subsection*{Q076}
Uma loja de roupas vende um certo modelo de jaqueta. O custo de aquisição de cada peça, pago ao fornecedor, é de R\$\,80,00. Uma pesquisa de mercado mostrou que, se a loja fixar o preço de venda em $p$ reais (com $80 < p < 150$), a quantidade de jaquetas vendidas por mês, em unidades, é dada por $q(p) = 300 - 2p$.

a) Escreva a expressão do lucro mensal $L(p)$ obtido com a venda das jaquetas, em função do preço $p$, sabendo que o lucro é o produto entre o lucro unitário (preço de venda menos custo de aquisição) e a quantidade vendida.

b) Determine o preço $p$ que maximiza o lucro mensal e calcule esse lucro máximo.

c) Confirme que esse preço realmente fornece o maior lucro possível, calculando o lucro para dois preços vizinhos — R\$\,114,00 e R\$\,116,00 — e comparando os resultados com o lucro obtido no preço ótimo encontrado no item (b).

d) A diretoria da loja quer saber para quais preços o lucro mensal ultrapassa R\$\,2.400,00. Determine esse intervalo de preços.

e) Suponha agora que o custo de aquisição da jaqueta suba de R\$\,80,00 para um valor genérico $c$ reais, mantendo-se a mesma relação de demanda $q(p) = 300-2p$. Escreva o lucro $L(p,c)$ em função de $p$ e $c$, obtenha o preço ótimo em função de $c$ e use essa expressão para determinar de quanto aumenta o preço ótimo se o custo subir de R\$\,80,00 para R\$\,90,00.

\textbf{Gabarito proposto:} a) $L(p) = -2p^2+460p-24000$; b) $p=115$ reais, lucro máximo $= 2450$ reais; c) $L(114)=L(116)=2448 < 2450$, confirmando o máximo; d) $110 < p < 120$; e) $p^*(c) = 75 + c/2$; aumento de R\$\,5,00 no preço ótimo quando o custo sobe R\$\,10,00.

\textbf{Habilidade declarada:} EM13MAT503 --- Investigar pontos de máximo ou de mínimo de funções quadráticas em contextos da Matemática Financeira ou da Cinemática, entre outros.
\end{samepage}

\noindent\rule{\linewidth}{0.4pt}
\textbf{Tem erro matemático?}\quad
\mbox{$\square$~não} \quad \mbox{$\square$~sim, aqui:} \underline{\hspace{5cm}}
\quad \mbox{$\square$~não sei dizer}

\textbf{Que nível de dificuldade esta questão tem para os seus alunos?}\quad
\mbox{$\square$~fácil} \quad \mbox{$\square$~média} \quad
\mbox{$\square$~difícil} \quad \mbox{$\square$~fora do alcance da turma}

\textbf{Você usaria esta questão?}\quad
\mbox{$\square$~aceita} \quad \mbox{$\square$~aceita com ajuste} \quad
\mbox{$\square$~recusada}

\textbf{Por quê?} \emph{(obrigatório quando não for ``aceita'')}

\underline{\hspace{\linewidth}}

\underline{\hspace{\linewidth}}
\vspace{0.6em}

\begin{samepage}
\subsection*{Q080}
Um laboratório de microbiologia monitora duas culturas de bactérias mantidas em condições ideais de nutrientes e temperatura, iniciando a contagem no instante $t=0$ (em horas).

\begin{itemize}[nosep,leftmargin=*]
\item A \textbf{Cultura A} parte de 500 bactérias e duplica sua população a cada 3 horas, sendo modelada por $P_A(t) = 500 \cdot 2^{t/3}$.
\item A \textbf{Cultura B} parte de 8000 bactérias e duplica sua população a cada 6 horas, sendo modelada por $P_B(t) = 8000 \cdot 2^{t/6}$.
\end{itemize}

Embora a Cultura B comece com uma população 16 vezes maior que a Cultura A, o pesquisador responsável afirma que, depois de certo tempo, a Cultura A ultrapassará a Cultura B e permanecerá com mais bactérias dali em diante, para sempre.

a) Determine o instante $t$ (em horas) em que as duas culturas apresentam exatamente a mesma quantidade de bactérias.

b) Sem calcular os valores das populações em instantes específicos, explique — comparando os tempos de duplicação (ou as taxas de crescimento) das duas culturas — por que, a partir do instante encontrado no item (a), a Cultura A permanecerá \textbf{para sempre} com mais bactérias que a Cultura B.

c) Confirme numericamente a afirmação do pesquisador calculando $P_A(30)$ e $P_B(30)$, isto é, as populações das duas culturas 30 horas após o início do experimento.

\textbf{Gabarito proposto:} a) $t = 24$ horas. b) Como a taxa de crescimento relativo de A ($\ln 2/3$) é maior que a de B ($\ln 2/6$), a razão $P_A(t)/P_B(t) = \frac{1}{16}\cdot 2^{t/6}$ é uma exponencial crescente que vale 1 em $t=24$h e cresce indefinidamente para $t>24$h; logo A ultrapassa B nesse instante e permanece maior para sempre. c) $P_A(30)=512000$ e $P_B(30)=256000$, confirmando que $P_A(30) > P_B(30)$.

\textbf{Habilidade declarada:} EM13MAT304 --- Resolver e elaborar problemas com funções exponenciais nos quais é necessário compreender e interpretar a variação das grandezas envolvidas, em contextos como o da Matemática Financeira e o do crescimento de seres vivos microscópicos, entre outros.
\end{samepage}

\noindent\rule{\linewidth}{0.4pt}
\textbf{Tem erro matemático?}\quad
\mbox{$\square$~não} \quad \mbox{$\square$~sim, aqui:} \underline{\hspace{5cm}}
\quad \mbox{$\square$~não sei dizer}

\textbf{Que nível de dificuldade esta questão tem para os seus alunos?}\quad
\mbox{$\square$~fácil} \quad \mbox{$\square$~média} \quad
\mbox{$\square$~difícil} \quad \mbox{$\square$~fora do alcance da turma}

\textbf{Você usaria esta questão?}\quad
\mbox{$\square$~aceita} \quad \mbox{$\square$~aceita com ajuste} \quad
\mbox{$\square$~recusada}

\textbf{Por quê?} \emph{(obrigatório quando não for ``aceita'')}

\underline{\hspace{\linewidth}}

\underline{\hspace{\linewidth}}
\vspace{0.6em}

\begin{samepage}
\subsection*{Q050}
Uma agência organiza passeios de van para grupos de turistas. A van comporta no máximo 60 passageiros, e a agência só realiza o passeio se houver pelo menos 15 inscritos.

O preço da passagem depende do tamanho do grupo: para 15 passageiros, o preço é R\$\,90,00 por pessoa. A cada passageiro adicional além de 15, a agência reduz o preço da passagem em R\$\,1,00 para TODOS os integrantes do grupo (não apenas para o excedente), como forma de incentivar grupos maiores.

Além da receita das passagens, a agência tem um custo de R\$\,500,00 fixo pelo aluguel da van, mais R\$\,20,00 por passageiro (referentes a seguro e lanche).

Determine o número de passageiros que a agência deve buscar para cada passeio, de modo a maximizar seu lucro, e calcule o valor desse lucro máximo.

\textbf{Gabarito proposto:} O lucro é máximo para 42 ou 43 passageiros, resultando em L = R\$\,1.306,00.

\textbf{Habilidade declarada:} EM13MAT503 --- Investigar pontos de máximo ou de mínimo de funções quadráticas em contextos da Matemática Financeira ou da Cinemática, entre outros.
\end{samepage}

\noindent\rule{\linewidth}{0.4pt}
\textbf{Tem erro matemático?}\quad
\mbox{$\square$~não} \quad \mbox{$\square$~sim, aqui:} \underline{\hspace{5cm}}
\quad \mbox{$\square$~não sei dizer}

\textbf{Que nível de dificuldade esta questão tem para os seus alunos?}\quad
\mbox{$\square$~fácil} \quad \mbox{$\square$~média} \quad
\mbox{$\square$~difícil} \quad \mbox{$\square$~fora do alcance da turma}

\textbf{Você usaria esta questão?}\quad
\mbox{$\square$~aceita} \quad \mbox{$\square$~aceita com ajuste} \quad
\mbox{$\square$~recusada}

\textbf{Por quê?} \emph{(obrigatório quando não for ``aceita'')}

\underline{\hspace{\linewidth}}

\underline{\hspace{\linewidth}}
\vspace{0.6em}

\begin{samepage}
\subsection*{Q036}
Uma equipe de engenharia testa, em túnel de vento, a força de resistência do ar (F, em newtons) que atua sobre uma peça, variando a velocidade do ar (v, em m/s) incidente sobre ela. Os resultados obtidos estão na tabela abaixo:

\begin{center}
\begin{tabular}{|p{0.172\linewidth}|p{0.172\linewidth}|p{0.172\linewidth}|p{0.172\linewidth}|p{0.172\linewidth}|}
\hline
\textbf{v (m/s)} & \textbf{1} & \textbf{2} & \textbf{3} & \textbf{4} \\
\hline
F (N) & 3 & 12 & 27 & 48 \\
\hline
\end{tabular}
\end{center}

Analisando os dados da tabela, qual expressão algébrica representa corretamente a força de resistência F em função da velocidade v?
\begin{enumerate}[label=(\alph*),nosep,leftmargin=*]
\item F = 3v²
\item F = 3v
\item F = 6v²
\item F = v² + 2
\end{enumerate}

\textbf{Gabarito proposto:} F(v) = 3v²

\textbf{Habilidade declarada:} EM13MAT502 --- Investigar relações entre números expressos em tabelas para representá-los no plano cartesiano, identificando padrões e criando conjecturas para generalizar e expressar algebricamente essa generalização, reconhecendo quando essa representação é de função polinomial de 2º grau do tipo y = ax².
\end{samepage}

\noindent\rule{\linewidth}{0.4pt}
\textbf{Tem erro matemático?}\quad
\mbox{$\square$~não} \quad \mbox{$\square$~sim, aqui:} \underline{\hspace{5cm}}
\quad \mbox{$\square$~não sei dizer}

\textbf{Que nível de dificuldade esta questão tem para os seus alunos?}\quad
\mbox{$\square$~fácil} \quad \mbox{$\square$~média} \quad
\mbox{$\square$~difícil} \quad \mbox{$\square$~fora do alcance da turma}

\textbf{Você usaria esta questão?}\quad
\mbox{$\square$~aceita} \quad \mbox{$\square$~aceita com ajuste} \quad
\mbox{$\square$~recusada}

\textbf{Por quê?} \emph{(obrigatório quando não for ``aceita'')}

\underline{\hspace{\linewidth}}

\underline{\hspace{\linewidth}}
\vspace{0.6em}

\begin{samepage}
\subsection*{Q034}
Duas empresas de aplicativo de transporte, A e B, calculam o valor da corrida em função da distância percorrida $x$ (em km). A empresa A cobra R\$\,2,50 por quilômetro rodado, sem nenhuma taxa adicional. A empresa B cobra uma taxa fixa de embarque de R\$\,3,00 mais R\$\,2,00 por quilômetro rodado. Sejam $y_A(x)$ e $y_B(x)$, em reais, os custos totais das corridas das empresas A e B, respectivamente, em função de $x$. Se essas duas funções forem representadas por retas em um plano cartesiano (custo $y$ no eixo vertical e distância $x$ no eixo horizontal), assinale a alternativa que descreve corretamente essas retas.
\begin{enumerate}[label=(\alph*),nosep,leftmargin=*]
\item A reta que representa a empresa A passa pela origem (0,0) e tem coeficiente angular 2,5, sendo uma função proporcional; a reta da empresa B corta o eixo y no ponto (0,3) e tem coeficiente angular 2, não sendo proporcional.
\item A reta que representa a empresa B passa pela origem (0,0), pois seu coeficiente angular é 2; a reta da empresa A corta o eixo y no ponto (0,3), pois seu coeficiente angular é 2,5.
\item Ambas as retas passam pela origem (0,0), já que representam custo por quilômetro rodado; a reta de A tem coeficiente angular 2,5 e a de B tem coeficiente angular 2.
\item A reta da empresa A corta o eixo y no ponto (0, 2,5); a reta da empresa B corta o eixo y no ponto (0,2) e tem coeficiente angular 3.
\end{enumerate}

\textbf{Gabarito proposto:} A reta que representa a empresa A passa pela origem $(0,0)$ e tem coeficiente angular 2,5 (função proporcional); a reta da empresa B corta o eixo y no ponto $(0,3)$ e tem coeficiente angular 2 (função afim, não proporcional).

\textbf{Habilidade declarada:} EM13MAT401 --- Converter representações algébricas de funções polinomiais de 1º grau para representações geométricas no plano cartesiano, distinguindo os casos nos quais o comportamento é proporcional, recorrendo ou não a softwares ou aplicativos de álgebra e geometria dinâmica.
\end{samepage}

\noindent\rule{\linewidth}{0.4pt}
\textbf{Tem erro matemático?}\quad
\mbox{$\square$~não} \quad \mbox{$\square$~sim, aqui:} \underline{\hspace{5cm}}
\quad \mbox{$\square$~não sei dizer}

\textbf{Que nível de dificuldade esta questão tem para os seus alunos?}\quad
\mbox{$\square$~fácil} \quad \mbox{$\square$~média} \quad
\mbox{$\square$~difícil} \quad \mbox{$\square$~fora do alcance da turma}

\textbf{Você usaria esta questão?}\quad
\mbox{$\square$~aceita} \quad \mbox{$\square$~aceita com ajuste} \quad
\mbox{$\square$~recusada}

\textbf{Por quê?} \emph{(obrigatório quando não for ``aceita'')}

\underline{\hspace{\linewidth}}

\underline{\hspace{\linewidth}}
\vspace{0.6em}

\newpage
\subsection*{Para terminar}

\textbf{1. Você usaria um sistema assim no seu planejamento? Por quê?}

\underline{\hspace{\linewidth}}

\underline{\hspace{\linewidth}}

\bigskip
\textbf{2. O que faltou nas questões que você viu?}

\underline{\hspace{\linewidth}}

\underline{\hspace{\linewidth}}

\bigskip
\textbf{3. O que atrapalhou --- na questão, no enunciado, no formato deste
material?}

\underline{\hspace{\linewidth}}

\underline{\hspace{\linewidth}}

\bigskip
Obrigado. Se quiser saber quais itens tinham erro e qual era, é só pedir.
\end{document}